\documentclass[11pt,a4paper]{article}
\usepackage[utf8]{inputenc}
\usepackage[T1]{fontenc}
\usepackage[brazilian]{babel}
\usepackage[margin=2cm]{geometry}
\usepackage{lmodern}
\usepackage{amssymb}
\usepackage{xcolor}
\usepackage{array}
\usepackage{booktabs}
\usepackage{enumitem}
\usepackage{titlesec}
\usepackage{fancyhdr}
\usepackage{parskip}

\definecolor{brandgreen}{HTML}{2E5E3E}
\definecolor{lightbg}{HTML}{F5F2EC}

\titleformat{\section}{\large\bfseries\color{brandgreen}}{}{0pt}{}[{\color{brandgreen}\titlerule}]
\titleformat{\subsection}{\bfseries\color{brandgreen}}{}{0pt}{}

\pagestyle{fancy}
\fancyhf{}
\fancyhead[L]{\small Querência ERP --- Avaliação de usabilidade}
\fancyhead[R]{\small Questionário SUS --- Painel administrativo}
\fancyfoot[C]{\small\thepage}
\renewcommand{\headrulewidth}{0.4pt}

\begin{document}

\begin{center}
  {\LARGE\bfseries\color{brandgreen} Questionário SUS}\\[1mm]
  {\large Painel Administrativo (Querência ERP) --- papel Diretoria}
\end{center}

\vspace{3mm}
\noindent\textbf{ID do participante} (código, não o nome): \makebox[4cm]{\hrulefill}

\vspace{2mm}
\noindent\textbf{Data:} \makebox[3.5cm]{\hrulefill}

\vspace{2mm}
\noindent\textbf{Entidade:} $\square$ CPF Pia do Sul \hspace{8mm} $\square$ CTG Sentinela da Querência

\vspace{4mm}
\noindent\colorbox{lightbg}{%
  \begin{minipage}{0.95\linewidth}
  \vspace{1mm}
  Pensando na sua experiência usando o \textbf{painel administrativo} agora há pouco, para cada
  frase abaixo marque o quanto você concorda ou discorda, numa escala de 1 (discordo totalmente) a
  5 (concordo totalmente). Não existe resposta certa ou errada.
  \vspace{1mm}
  \end{minipage}%
}

\vspace{4mm}
\renewcommand{\arraystretch}{1.6}
\noindent\begin{tabular}{@{}p{0.6cm}p{8.2cm}p{1.0cm}p{1.0cm}p{1.0cm}p{1.0cm}p{1.0cm}@{}}
  \toprule
   & \textbf{Afirmação} & \textbf{1} & \textbf{2} & \textbf{3} & \textbf{4} & \textbf{5} \\
  \midrule
  1 & Eu acho que gostaria de usar o painel administrativo com frequência. & $\bigcirc$ & $\bigcirc$ & $\bigcirc$ & $\bigcirc$ & $\bigcirc$ \\
  2 & Eu achei o painel administrativo complexo sem necessidade. & $\bigcirc$ & $\bigcirc$ & $\bigcirc$ & $\bigcirc$ & $\bigcirc$ \\
  3 & Eu achei o painel administrativo fácil de usar. & $\bigcirc$ & $\bigcirc$ & $\bigcirc$ & $\bigcirc$ & $\bigcirc$ \\
  4 & Eu acho que precisaria da ajuda de uma pessoa com conhecimentos técnicos para usar o painel administrativo. & $\bigcirc$ & $\bigcirc$ & $\bigcirc$ & $\bigcirc$ & $\bigcirc$ \\
  5 & Eu achei que as várias funções do painel administrativo estavam bem integradas. & $\bigcirc$ & $\bigcirc$ & $\bigcirc$ & $\bigcirc$ & $\bigcirc$ \\
  6 & Eu achei que houve muita inconsistência no painel administrativo. & $\bigcirc$ & $\bigcirc$ & $\bigcirc$ & $\bigcirc$ & $\bigcirc$ \\
  7 & Eu imagino que a maioria das pessoas aprenderia a usar o painel administrativo rapidamente. & $\bigcirc$ & $\bigcirc$ & $\bigcirc$ & $\bigcirc$ & $\bigcirc$ \\
  8 & Eu achei o painel administrativo muito complicado de usar. & $\bigcirc$ & $\bigcirc$ & $\bigcirc$ & $\bigcirc$ & $\bigcirc$ \\
  9 & Eu me senti confiante usando o painel administrativo. & $\bigcirc$ & $\bigcirc$ & $\bigcirc$ & $\bigcirc$ & $\bigcirc$ \\
  10 & Eu precisei aprender várias coisas novas antes de conseguir usar o painel administrativo. & $\bigcirc$ & $\bigcirc$ & $\bigcirc$ & $\bigcirc$ & $\bigcirc$ \\
  \bottomrule
\end{tabular}

\vspace{5mm}
\noindent\textbf{Algum comentário livre sobre a experiência?} (opcional)

\vspace{2mm}
\noindent\hrulefill

\vspace{2mm}
\noindent\hrulefill

\vfill
\noindent\footnotesize\textit{Nota de apuração (não faz parte do formulário entregue ao
participante):} itens ímpares (1,3,5,7,9) contribuem com (valor $-$ 1); itens pares (2,4,6,8,10)
com (5 $-$ valor); a soma das dez contribuições (0--40) é multiplicada por 2,5, resultando numa
pontuação de 0 a 100. Faixas de interpretação (Bangor, Kortum e Miller, 2009): $\geq$ 80,3
excelente; 68--80,2 bom; 51--67,9 ok/aceitável; $<$ 51 ruim. Reportar separadamente da pontuação do
app, usando a \textbf{mediana} (não a média) ao agregar por papel. Instrumento de Brooke (1996), na
adaptação PT-BR de Lourenço, Carmona e Lopes (2022).

\end{document}
